\documentclass[conference]{IEEEtran}

\usepackage{amsmath}
\usepackage{amssymb}
\usepackage{amsthm}
\usepackage{cite}
\usepackage{algorithm}
\usepackage{algorithmic}
\usepackage{booktabs}

\allowdisplaybreaks

% ----- Standard theorem environments -----
\newtheorem{proposition}{Proposition}
\newtheorem{theorem}{Theorem}
\newtheorem{lemma}{Lemma}
\newtheorem{corollary}{Corollary}
\newtheorem{remark}{Remark}
\newtheorem{definition}{Definition}

\begin{document}

\title{Mathematical Foundations of the WCS-GNCCP\\Molecular Graph Kernel with Hierarchical\\2D--3D Alignment}

\author{%
\IEEEauthorblockN{Allampati Chenchu Sanhitha Reddy,
Shaik Abdus Sattar,
Geethika G Nair,
Meera J Vaishnav}
\IEEEauthorblockA{Department of Computer Science and Engineering\\
Amrita Vishwa Vidyapeetham, Coimbatore\\
Guide: Dr.\ T Ramraj}}

\maketitle

\begin{abstract}
We present the mathematical foundations of a hierarchical molecular
similarity kernel that combines 2D topological matching with 3D
geometric alignment.  The 2D stage solves a Weighted Common Subgraph
(WCS) problem via Graduated Non-Convexity and Concavity Procedure
(GNCCP) with Frank-Wolfe optimization, producing both a similarity
score and a bijective atom correspondence.  We derive a closed-form
structural gradient (Theorem~\ref{thm:gradient}), prove a rank-one
factorization that eliminates one cubic sub-operation per gradient
evaluation (Lemma~\ref{lem:U-outer}), and establish scale invariance
via Frobenius normalization (Proposition~\ref{prop:normalization}).
The bijective map produced by Stage~1 reduces the 3D alignment
(Stage~2) to a single $O(n)$ Kabsch rotation, yielding a hybrid kernel
with per-pair complexity dominated by the 2D stage at
$O(T_\zeta T_{\mathrm{FW}} n^3)$ --- empirically competitive with
Random Walk and Optimal Assignment kernels on drug-like molecules.
\end{abstract}


% ===================================================================
\section{Problem Formulation}
\label{sec:formulation}
% ===================================================================

Let $G$ and $H$ be two molecular graphs with $M$ and $N$ atoms, respectively.
Their topology is encoded by weighted adjacency matrices
$\mathbf{A}_G \!\in\! \mathbb{R}^{M \times M}$ and
$\mathbf{A}_H \!\in\! \mathbb{R}^{N \times N}$,
where each entry reflects bond type (single, double, triple, aromatic)
and topology (ring membership, conjugation).
A cost matrix $\mathbf{C} \!\in\! \mathbb{R}^{M \times N}$ encodes pairwise
atom-feature dissimilarity.

\begin{definition}[Partial Permutation Polytope]
\label{def:polytope}
For an integer $L \leq \min(M, N)$, the partial permutation polytope is
\begin{align}
\mathcal{P}_L = \big\{ \mathbf{X} \in \{0,1\}^{M \times N} \;\big|\;
  &\mathbf{X}\mathbf{1}_N \leq \mathbf{1}_M,\notag\\
  &\mathbf{X}^{T}\mathbf{1}_M \leq \mathbf{1}_N,\notag\\
  &\mathbf{1}_M^{T}\mathbf{X}\mathbf{1}_N = L \big\},
\label{eq:polytope}
\end{align}
and its convex relaxation $\mathcal{D}_L$ replaces $\{0,1\}^{M \times N}$ with $[0,1]^{M \times N}$.
An entry $X_{ij} = 1$ indicates that atom~$i$ in~$G$ is matched to atom~$j$ in~$H$.
\end{definition}

\begin{definition}[WCS Objective]
\label{def:wcs-objective}
The Weighted Common Subgraph (WCS) objective is
\begin{equation}
F(\mathbf{X}) = \alpha\, f(\mathbf{X})
  + (1 - \alpha)\operatorname{tr}\!\big(\mathbf{C}^{T}\!\mathbf{X}\big),
\label{eq:wcs-objective}
\end{equation}
where $\alpha \!\in\! [0,1]$ balances structure against atom features,
\begin{equation}
f(\mathbf{X})
  = \big\| \mathbf{U} \circ \mathbf{A}_G
    - \mathbf{X}\mathbf{A}_H\mathbf{X}^{T} \big\|_F^2
\label{eq:structural}
\end{equation}
is the structural mismatch term, $\circ$ denotes the Hadamard product, and
\begin{equation}
\mathbf{U}
  = \mathbf{X}\,\mathbf{1}_{N \times N}\,\mathbf{X}^{T}
  \in \mathbb{R}^{M \times M}
\label{eq:U-def}
\end{equation}
is the co-matching indicator.
\end{definition}


% ===================================================================
\section{Computational Identity for $\mathbf{U}$}
\label{sec:U-identity}
% ===================================================================

The matrix $\mathbf{U}$ in~\eqref{eq:U-def} involves the
$N \!\times\! N$ all-ones matrix, requiring $O(MN^2 + M^2 N)$ work.
We show that $\mathbf{U}$ admits a rank-one factorization
that reduces this to $O(MN + M^2)$.

\begin{lemma}[Rank-One Factorization of~$\mathbf{U}$]
\label{lem:U-outer}
Let\, $\mathbf{r} = \mathbf{X}\mathbf{1}_N \!\in\! \mathbb{R}^{M}$ be the
row-sum vector. Then
\begin{equation}
\mathbf{U}
  = \mathbf{X}\,\mathbf{1}_{N \times N}\,\mathbf{X}^{T}
  = \mathbf{r}\,\mathbf{r}^{T}.
\label{eq:U-outer}
\end{equation}
\end{lemma}

\begin{proof}
Factor the all-ones matrix as
$\mathbf{1}_{N \times N} = \mathbf{1}_N\,\mathbf{1}_N^{T}$:
\begin{align}
\mathbf{U}
  &= \mathbf{X}\bigl(\mathbf{1}_N\,\mathbf{1}_N^{T}\bigr)\mathbf{X}^{T}
\notag\\
  &= \bigl(\mathbf{X}\mathbf{1}_N\bigr)
     \bigl(\mathbf{X}\mathbf{1}_N\bigr)^{T}
  = \mathbf{r}\,\mathbf{r}^{T}.\label{eq:U-proof}
\end{align}
\end{proof}

\begin{remark}
$\mathbf{U}$ is symmetric PSD with rank one.
For binary $\mathbf{X} \!\in\! \mathcal{P}_L$ each $r_i \!\in\! \{0,1\}$,
so $U_{ik} = 1$ iff both atoms $i,k$ are matched.
\end{remark}


% ===================================================================
\section{Gradient of the Structural Term}
\label{sec:gradient}
% ===================================================================

We derive $\nabla_{\!\mathbf{X}} f$ for the structural cost~\eqref{eq:structural}.

\begin{theorem}[Structural Gradient]
\label{thm:gradient}
Define the residual
$\mathbf{R} = \mathbf{U} \circ \mathbf{A}_G
              - \mathbf{X}\mathbf{A}_H\mathbf{X}^{T}$.
Then
\begin{equation}
\nabla_{\!\mathbf{X}} f
  = 4\bigl(\mathbf{R} \circ \mathbf{A}_G\bigr)
    \mathbf{X}\,\mathbf{1}_{N \times N}
  - 4\,\mathbf{R}\,\mathbf{X}\,\mathbf{A}_H.
\label{eq:gradient}
\end{equation}

Moreover, using Lemma~\ref{lem:U-outer}:
\begin{equation}
\bigl(\mathbf{R} \circ \mathbf{A}_G\bigr)
  \mathbf{X}\,\mathbf{1}_{N \times N}
= \Bigl[\bigl(\mathbf{R} \circ \mathbf{A}_G\bigr)\mathbf{X}\Bigr]
  \mathbf{1}_N\mathbf{1}_N^{T},
\label{eq:gradient-broadcast}
\end{equation}
a rank-one matrix whose $i$-th row repeats
$\sum_{j}\bigl[(\mathbf{R}\!\circ\!\mathbf{A}_G)\mathbf{X}\bigr]_{ij}$.
\end{theorem}

\begin{proof}
Write $f = \operatorname{tr}(\mathbf{R}^{T}\mathbf{R})$.
By the chain rule,
\begin{equation}
\mathrm{d}f = 2\operatorname{tr}\!\bigl(\mathbf{R}^{T}\,\mathrm{d}\mathbf{R}\bigr).
\label{eq:df}
\end{equation}

\emph{Part~1 (Hadamard term).}
Since $\mathbf{A}_G$ is constant,
\begin{equation}
\mathrm{d}\bigl(\mathbf{U} \circ \mathbf{A}_G\bigr)
  = \bigl(\mathrm{d}\mathbf{U}\bigr) \circ \mathbf{A}_G.
\label{eq:dUA}
\end{equation}
From $\mathbf{U} = \mathbf{r}\,\mathbf{r}^{T}$
with $\mathbf{r} = \mathbf{X}\mathbf{1}_N$:
\begin{equation}
\mathrm{d}\mathbf{U}
  = (\mathrm{d}\mathbf{X}\,\mathbf{1}_N)\,\mathbf{r}^{T}
  + \mathbf{r}\,(\mathrm{d}\mathbf{X}\,\mathbf{1}_N)^{T}.
\label{eq:dU}
\end{equation}

\emph{Part~2 (Bilinear term).}
By the product rule,
\begin{align}
\mathrm{d}\bigl(\mathbf{X}\mathbf{A}_H\mathbf{X}^{T}\bigr)
  &= (\mathrm{d}\mathbf{X})\,\mathbf{A}_H\,\mathbf{X}^{T}
\notag\\
  &\quad+ \mathbf{X}\,\mathbf{A}_H\,(\mathrm{d}\mathbf{X})^{T}.
\label{eq:dXAX}
\end{align}

\emph{Part~3 (Combining).}
Substitute \eqref{eq:dUA}--\eqref{eq:dXAX} into~\eqref{eq:df}
and apply the Hadamard trace identity
\begin{equation*}
\operatorname{tr}\!\bigl(\mathbf{A}^{T}(\mathbf{B}\!\circ\!\mathbf{C})\bigr)
= \operatorname{tr}\!\bigl((\mathbf{A}\!\circ\!\mathbf{B})^{T}\mathbf{C}\bigr).
\end{equation*}
Collecting terms in $\mathrm{d}\mathbf{X}$: the Hadamard contribution
yields $(\mathbf{R}\!\circ\!\mathbf{A}_G)\mathbf{X}\mathbf{1}_{N\times N}$
and the bilinear contribution yields $-\mathbf{R}\mathbf{X}\mathbf{A}_H$.
Since $\mathbf{R}$, $\mathbf{A}_G$, $\mathbf{U}$, and
$\mathbf{X}\mathbf{A}_H\mathbf{X}^{T}$ are all symmetric,
each of the two terms in~\eqref{eq:dU} and~\eqref{eq:dXAX}
contributes equally, producing the factor of~$4$.
\end{proof}

\begin{remark}
A common error is to differentiate $\mathbf{U} \circ \mathbf{A}_G$
as if it were the matrix product $\mathbf{U}\mathbf{A}_G$.
The Hadamard derivative is
$\mathrm{d}(\mathbf{U}\!\circ\!\mathbf{A}_G) = (\mathrm{d}\mathbf{U})\!\circ\!\mathbf{A}_G$
(element-wise), \emph{not} $(\mathrm{d}\mathbf{U})\mathbf{A}_G$.
\end{remark}


% ===================================================================
\section{The GNCCP Continuation Method}
\label{sec:gnccp}
% ===================================================================

Direct minimization of $F$ over $\mathcal{D}_L$ is nonconvex.
GNCCP introduces $\zeta \!\in\! [-1,1]$ to deform the problem
from convex to concave.

\begin{definition}[GNCCP Objective]
\label{def:Jzeta}
For $\zeta \!\in\! [-1, 1]$, define
\begin{equation}
J_\zeta(\mathbf{X}) =
\begin{cases}
(1 - \zeta)\,F(\mathbf{X}) + \zeta\,\|\mathbf{X}\|_F^2
  & \zeta \geq 0, \\[6pt]
(1 + \zeta)\,F(\mathbf{X}) + \zeta\,\|\mathbf{X}\|_F^2
  & \zeta < 0.
\end{cases}
\label{eq:Jzeta}
\end{equation}
\end{definition}

\begin{proposition}[Endpoint Properties]
\label{prop:endpoints}\hfill
\begin{enumerate}
\item $\zeta \!=\! +1$:\;
  $J_{+1} = \|\mathbf{X}\|_F^2$ is strictly convex.
  Its unique minimizer over $\mathcal{D}_L$ is the
  uniform matrix $\mathbf{X}^{(0)} = \tfrac{L}{MN}\,\mathbf{1}_{M\times N}$.

\item $\zeta \!=\! 0$:\;
  $J_0 = F(\mathbf{X})$, the original WCS objective.

\item $\zeta \!=\! -1$:\;
  $J_{-1} = -\|\mathbf{X}\|_F^2$ is strictly concave.
  Its minimizers over $\mathcal{D}_L$ are vertices,
  i.e., partial permutation matrices in $\mathcal{P}_L$.
\end{enumerate}
\end{proposition}

\begin{proof}
(1)~$\|\mathbf{X}\|_F^2 = \sum_{ij} X_{ij}^2$ is strictly convex.
The unique minimizer distributes mass uniformly:
$X_{ij} = L/(MN)$ satisfies all constraints
and minimizes the sum of squares by Schur convexity.

(2)~Direct substitution.

(3)~$-\|\mathbf{X}\|_F^2$ is strictly concave.
A strictly concave function attains its minimum
over a compact convex polytope at a vertex.
The vertices of~$\mathcal{D}_L$ are
exactly the elements of~$\mathcal{P}_L$.
\end{proof}

\begin{proposition}[Gradient of $J_\zeta$]
\label{prop:grad-Jzeta}
\begin{equation}
\nabla_{\!\mathbf{X}} J_\zeta\!=\!
\begin{cases}
(1 - \zeta)\,\nabla_{\!\mathbf{X}} F + 2\zeta\,\mathbf{X}
  & \zeta \geq 0, \\[4pt]
(1 + \zeta)\,\nabla_{\!\mathbf{X}} F + 2\zeta\,\mathbf{X}
  & \zeta < 0,
\end{cases}
\label{eq:grad-Jzeta}
\end{equation}
where\;
$\nabla_{\!\mathbf{X}} F
  = \alpha\,\nabla_{\!\mathbf{X}} f + (1\!-\!\alpha)\,\mathbf{C}$\;
and $\nabla_{\!\mathbf{X}} f$ is given by
Theorem~\ref{thm:gradient}.
\end{proposition}

\begin{proof}
$\nabla\|\mathbf{X}\|_F^2 = 2\mathbf{X}$.
Apply linearity to~\eqref{eq:Jzeta}.
\end{proof}


% ===================================================================
\section{Frank-Wolfe Optimization}
\label{sec:frank-wolfe}
% ===================================================================

For each fixed~$\zeta$ we solve
$\min_{\mathbf{X} \in \mathcal{D}_L} J_\zeta(\mathbf{X})$
via Frank-Wolfe (conditional gradient).

\begin{definition}[Linear Minimization Oracle]
\label{def:lmo}
At iterate $\mathbf{X}_t$, the FW direction is
\begin{equation}
\mathbf{Y}_t
  = \arg\min_{\mathbf{Y} \in \mathcal{D}_L}
    \bigl\langle \nabla J_\zeta(\mathbf{X}_t),\,\mathbf{Y}\bigr\rangle.
\label{eq:lmo}
\end{equation}
Since the minimum of a linear function over a polytope
is attained at a vertex, $\mathbf{Y}_t \!\in\! \mathcal{P}_L$.
The optimal vertex is found by the Hungarian algorithm
in $O\!\bigl(\min(M,N)^3\bigr)$ time.
\end{definition}

\begin{theorem}[Frank-Wolfe Duality Gap]
\label{thm:fw-gap}
Let $J_\zeta$ be convex on $\mathcal{D}_L$ and let
$\mathbf{X}^{\!*} = \arg\min_{\mathcal{D}_L} J_\zeta$.
The Frank-Wolfe gap
\begin{equation}
g_t
  = \bigl\langle \nabla J_\zeta(\mathbf{X}_t),\;
    \mathbf{X}_t - \mathbf{Y}_t \bigr\rangle
\label{eq:fw-gap}
\end{equation}
satisfies:\;
\textup{(i)}~$g_t \geq 0$;\;
\textup{(ii)}~$g_t \geq J_\zeta(\mathbf{X}_t) - J_\zeta(\mathbf{X}^{\!*})$;\;
\textup{(iii)}~$g_t = 0 \Leftrightarrow \mathbf{X}_t$ is optimal.
\end{theorem}

\begin{proof}
\emph{(i)}
$\mathbf{Y}_t$ minimizes the linear form over $\mathcal{D}_L$
and $\mathbf{X}_t \!\in\! \mathcal{D}_L$, so
$\langle\nabla J_\zeta,\mathbf{Y}_t\rangle
 \leq \langle\nabla J_\zeta,\mathbf{X}_t\rangle$,
giving $g_t \geq 0$.

\emph{(ii)}
By convexity:
\begin{equation}
J_\zeta(\mathbf{X}^{\!*})
  \geq J_\zeta(\mathbf{X}_t)
  + \bigl\langle\nabla J_\zeta(\mathbf{X}_t),\,
    \mathbf{X}^{\!*}\!-\!\mathbf{X}_t\bigr\rangle.
\label{eq:convexity-lb}
\end{equation}
Since $\mathbf{Y}_t$ is the minimizer and
$\mathbf{X}^{\!*}\!\in\!\mathcal{D}_L$:
\begin{equation}
\bigl\langle\nabla J_\zeta(\mathbf{X}_t),\,
  \mathbf{X}^{\!*}\!-\!\mathbf{X}_t\bigr\rangle
\geq \bigl\langle\nabla J_\zeta(\mathbf{X}_t),\,
  \mathbf{Y}_t\!-\!\mathbf{X}_t\bigr\rangle
= -g_t.
\label{eq:gap-ineq}
\end{equation}
Combining \eqref{eq:convexity-lb}--\eqref{eq:gap-ineq}:
$g_t \geq J_\zeta(\mathbf{X}_t) - J_\zeta(\mathbf{X}^{\!*})$.

\emph{(iii)}
If $\mathbf{X}_t$ is optimal, first-order optimality gives
$\langle\nabla J_\zeta,\mathbf{Y}\!-\!\mathbf{X}_t\rangle \!\geq\! 0$
$\forall\,\mathbf{Y}\!\in\!\mathcal{D}_L$, so $g_t = 0$.
Conversely, $g_t = 0$ with~(ii) gives
$J_\zeta(\mathbf{X}_t) \leq J_\zeta(\mathbf{X}^{\!*})$.
\end{proof}

\begin{remark}
When $J_\zeta$ is nonconvex ($\zeta \!<\! 0$),
the gap~\eqref{eq:fw-gap} is no longer a certified
suboptimality bound, but $g_t = 0$ still implies
$\mathbf{X}_t$ is a stationary point.
We use $g_t < \varepsilon$ as early termination for all~$\zeta$.
\end{remark}


% ===================================================================
\section{Kernel Properties}
\label{sec:kernel}
% ===================================================================

\begin{proposition}[Exponential Graph Similarity Kernel]
\label{prop:psd}
The exponential mapping
\begin{equation}
k(G, H) = \exp\!\Bigl(-\frac{\gamma\, F^{\!*}_{GH}}{L}\Bigr)
\label{eq:kernel}
\end{equation}
where $\gamma > 0$ is a bandwidth parameter and $F^{\!*}_{GH}$
is the optimal WCS objective, defines a discriminative
similarity measure satisfying:
\begin{enumerate}
\item $k(G, H) \in (0, 1]$ for all graph pairs;
\item $k(G, G) = 1$ when $G$ is matched to itself;
\item $k$ is monotonically decreasing in $F^{\!*}_{GH}$.
\end{enumerate}
\end{proposition}

\begin{remark}
The kernel~\eqref{eq:kernel} is generally \emph{not}
positive semi-definite in the Mercer sense,
since optimal-assignment and graph-edit distances do not
satisfy conditional negative definiteness
(cf.\ Schoenberg~\cite{schoenberg1938metric}).
This is not a limitation for distance-based classifiers
such as $k$-NN (which rely on ranking, not inner products)
or for SVM with eigenvalue ridge correction
($\mathbf{K}_{\text{psd}} = \mathbf{K} + \mu\mathbf{I}$,
$\mu \geq |\lambda_{\min}|$).
\end{remark}

\begin{remark}
$\gamma$ controls bandwidth:
small $\gamma$ spreads kernel values (all pairs appear similar),
large $\gamma$ concentrates mass on near-identical pairs.
We default to $\gamma = 0.5$ and normalize by $\max(L,1)$
for size invariance.
\end{remark}


% ===================================================================
\section{Scale Invariance}
\label{sec:normalization}
% ===================================================================

\begin{proposition}[Frobenius Normalization]
\label{prop:normalization}
Let $\hat{\mathbf{A}}_G = \mathbf{A}_G / \|\mathbf{A}_G\|_F$
and $\hat{\mathbf{A}}_H = \mathbf{A}_H / \|\mathbf{A}_H\|_F$.
Then
\begin{equation}
\hat{f}(\mathbf{X})
  = \bigl\| \mathbf{U} \circ \hat{\mathbf{A}}_G
    - \mathbf{X}\hat{\mathbf{A}}_H\mathbf{X}^{T} \bigr\|_F^2
  \leq 2
\label{eq:normalized-cost}
\end{equation}
for all $\mathbf{X} \!\in\! \mathcal{D}_L$,
independently of the number of bonds.
\end{proposition}

\begin{proof}
By the identity
$\|\mathbf{A}\!-\!\mathbf{B}\|_F^2
= \|\mathbf{A}\|_F^2 + \|\mathbf{B}\|_F^2
- 2\langle\mathbf{A},\mathbf{B}\rangle$,
it suffices to bound each squared norm.

For the Hadamard term, using
$\|\mathbf{A}\!\circ\!\mathbf{B}\|_F
\leq \|\mathbf{A}\|_{\max}\|\mathbf{B}\|_F$
with $U_{ik} = r_i r_k \leq 1$ for
$\mathbf{X}\!\in\!\mathcal{D}_L$:
\begin{equation}
\bigl\|\mathbf{U}\circ\hat{\mathbf{A}}_G\bigr\|_F
  \leq 1 \cdot \|\hat{\mathbf{A}}_G\|_F = 1.
\label{eq:bound1}
\end{equation}

For the bilinear term, since $\|\mathbf{X}\|_2 \leq 1$
for sub-stochastic~$\mathbf{X}$:
\begin{equation}
\bigl\|\mathbf{X}\hat{\mathbf{A}}_H\mathbf{X}^{T}\bigr\|_F
  \leq \|\mathbf{X}\|_2^2\,\|\hat{\mathbf{A}}_H\|_F
  \leq 1.
\label{eq:bound2}
\end{equation}

Therefore
$\hat{f} \leq \|\mathbf{U}\!\circ\!\hat{\mathbf{A}}_G\|_F^2
  + \|\mathbf{X}\hat{\mathbf{A}}_H\mathbf{X}^{T}\|_F^2
  \leq 2$.
\end{proof}

\begin{remark}
Without normalization, $f(\mathbf{X})$ scales with
$\|\mathbf{A}_G\|_F^2\!+\!\|\mathbf{A}_H\|_F^2$,
which grows linearly in the number of bonds.
For~30 bonds at weight~${\sim}1.5$,
$\|\mathbf{A}_G\|_F \!\approx\! 9.5$,
making the structural cost ${\sim}90\!\times$ larger
than the appearance cost
$\operatorname{tr}(\mathbf{C}^{T}\!\mathbf{X}) \!\in\! [0,L]$.
Normalization restores the balance intended by~$\alpha$.
\end{remark}


% ===================================================================
\section{Projection onto $\mathcal{D}_L$}
\label{sec:projection}
% ===================================================================

\begin{proposition}[Sinkhorn-Knopp Projection]
\label{prop:sinkhorn}
Let $\mathbf{X}\!\in\!\mathbb{R}_{\geq 0}^{M \times N}$ with
$\|\mathbf{X}\|_1 > 0$. The iteration

\smallskip\noindent
\textup{1.}~\textbf{Scale}: $\mathbf{X} \leftarrow \mathbf{X}\cdot(L/\|\mathbf{X}\|_1)$.\\
\textup{2.}~\textbf{Row cap}: if $(\mathbf{X}\mathbf{1}_N)_i > 1$,
  divide row~$i$ by $(\mathbf{X}\mathbf{1}_N)_i$.\\
\textup{3.}~\textbf{Col cap}: if $(\mathbf{X}^{T}\!\mathbf{1}_M)_j > 1$,
  divide col~$j$ by $(\mathbf{X}^{T}\!\mathbf{1}_M)_j$.\\
\textup{4.}~Repeat \textup{1--3} until convergence.

\smallskip\noindent
converges to a feasible point in $\mathcal{D}_L$
when $L \leq \min(M,N)$.
\end{proposition}

\begin{proof}[Proof sketch]
Steps~2--3 are alternating projections onto
$\mathcal{R}\!=\!\{\mathbf{X}\!\geq\!0 : \mathbf{X}\mathbf{1}\!\leq\!\mathbf{1}\}$
and
$\mathcal{C}\!=\!\{\mathbf{X}\!\geq\!0 : \mathbf{X}^{T}\!\mathbf{1}\!\leq\!\mathbf{1}\}$,
both convex and closed.
Step~1 enforces $\|\mathbf{X}\|_1\!=\!L$.
By Dykstra's theorem, this converges when the intersection
is non-empty.
The uniform matrix $\frac{L}{MN}\mathbf{1}_{M\times N}$
lies in all three constraint sets, so the intersection is non-empty.
\end{proof}


% ===================================================================
\section{Computational Complexity Analysis}
\label{sec:complexity}
% ===================================================================

We provide an honest, operation-level complexity analysis
of each Frank-Wolfe iteration.

\begin{theorem}[Per-Iteration Complexity]
\label{thm:complexity}
Let $M,N$ denote the atom counts and let $d$ be the maximum
vertex degree of the molecular graphs.
One Frank-Wolfe iteration of Algorithm~\ref{alg:gnccp-fw}
has complexity
\begin{equation}
\Theta\!\bigl(\underbrace{M^2 N}_{\mathbf{X}\mathbf{A}_H\mathbf{X}^T}
+ \underbrace{\min(M,N)^3}_{\text{Hungarian}}\bigr).
\label{eq:iter-complexity}
\end{equation}
For $M \!\approx\! N$, this simplifies to $\Theta(n^3)$.
\end{theorem}

\begin{proof}
We enumerate the dominant operations per iteration:

\smallskip\noindent
\emph{(a) Row sums:}
$\mathbf{r} = \mathbf{X}\mathbf{1}_N$ costs $O(MN)$.

\smallskip\noindent
\emph{(b) Co-matching indicator (Lemma~\ref{lem:U-outer}):}
$\mathbf{U} = \mathbf{r}\mathbf{r}^T$ costs $O(M^2)$.
Without Lemma~\ref{lem:U-outer}, the na\"ive computation
via~\eqref{eq:U-def} costs $O(MN^2 + M^2 N)$.

\smallskip\noindent
\emph{(c) Bilinear form:}
$\mathbf{P} = \mathbf{X}\mathbf{A}_H$ costs $O(M \cdot \mathrm{nnz}(\mathbf{A}_H))$;
for molecular graphs with bounded degree~$d$,
$\mathrm{nnz}(\mathbf{A}_H) = O(dN)$,
so $\mathbf{P}$ costs $O(MdN) \approx O(MN)$.
However, $\mathbf{P}\mathbf{X}^T$ is an
$(M\!\times\!N)(N\!\times\!M)$ product costing $O(M^2 N)$.
This is irreducible for dense~$\mathbf{X}$.

\smallskip\noindent
\emph{(d) Gradient terms:}
$({\mathbf{R}\!\circ\!\mathbf{A}_G})\mathbf{X}$ costs
$O(\mathrm{nnz}(\mathbf{A}_G)\!\cdot\!N) = O(dMN)$;
$\mathbf{R}\mathbf{X}\mathbf{A}_H$ costs $O(M^2 N + MdN)$.

\smallskip\noindent
\emph{(e) Hungarian algorithm:}
$O(\min(M,N)^3)$; this is tight (Kuhn--Munkres).

\smallskip\noindent
\emph{(f) Projection:}
$O(MN)$ per Sinkhorn iteration, bounded at 20 iterations.

\smallskip\noindent
The dominant terms are~(c) and~(e), giving~\eqref{eq:iter-complexity}.
\end{proof}

\begin{corollary}[Reduction from Lemma~\ref{lem:U-outer}]
\label{cor:reduction}
The rank-one factorization eliminates one $O(MN^2 + M^2 N)$
operation per gradient evaluation, reducing it to $O(MN + M^2)$.
For $M \!=\! N \!=\! n$, this is a reduction from $O(n^3)$ to $O(n^2)$
for the $\mathbf{U}$ computation alone.
The overall per-iteration complexity remains
$\Theta(n^3)$ due to the bilinear form~(c) and
Hungarian assignment~(e).
\end{corollary}

\begin{corollary}[Total Per-Pair Complexity]
\label{cor:total}
With $T_\zeta$ GNCCP outer iterations and $T_{\mathrm{FW}}$
Frank-Wolfe inner iterations, the total per-pair cost is
\begin{equation}
O\!\bigl(T_\zeta\, T_{\mathrm{FW}}\, n^3\bigr).
\label{eq:total-complexity}
\end{equation}
In practice, early termination via the FW duality gap
(Theorem~\ref{thm:fw-gap}) reduces the effective
iteration count well below the worst case.
\end{corollary}

\begin{remark}[Practical Competitiveness]
For drug-like molecules ($n \!\leq\! 50$),
the cubic cost per iteration is sub-millisecond.
With $(T_\zeta, T_{\mathrm{FW}}) = (4, 4)$ and hydrogen
removal ($n \!\approx\! 30$ heavy atoms), the measured
per-pair time is ${\sim}6$\,ms on CPU ---
competitive with Random Walk ($O(n^3)$, ${\sim}10$--$30$\,ms)
and Optimal Assignment kernels ($O(n^3)$, ${\sim}10$--$50$\,ms),
while producing a richer output (bijective atom map
in addition to similarity score).
\end{remark}


% ===================================================================
\section{Stage 2: 3D Alignment via Kabsch Rotation}
\label{sec:3d-alignment}
% ===================================================================

The bijective atom correspondence $\mathbf{X}^*$ from Stage~1
induces a set of $L$ matched atom pairs
$\mathcal{M} = \{(i,j) : X^*_{ij} = 1\}$.
We exploit this mapping to perform 3D alignment in $O(L)$ time.

\begin{definition}[Matched Coordinate Sets]
\label{def:matched-coords}
Let $\mathbf{p}_1, \ldots, \mathbf{p}_L \in \mathbb{R}^3$
and $\mathbf{q}_1, \ldots, \mathbf{q}_L \in \mathbb{R}^3$
be the 3D coordinates of the matched atoms in molecules
$G$ and $H$ respectively, ordered according to $\mathcal{M}$.
Define centered coordinates
$\bar{\mathbf{p}}_l = \mathbf{p}_l - \boldsymbol{\mu}_P$ and
$\bar{\mathbf{q}}_l = \mathbf{q}_l - \boldsymbol{\mu}_Q$,
where $\boldsymbol{\mu}_P = \frac{1}{L}\sum_l \mathbf{p}_l$
and $\boldsymbol{\mu}_Q = \frac{1}{L}\sum_l \mathbf{q}_l$.
\end{definition}

\begin{theorem}[Kabsch Alignment]
\label{thm:kabsch}
The rotation matrix $\mathbf{R}^* \in \mathrm{SO}(3)$ that
minimizes the RMSD
\begin{equation}
\mathrm{RMSD}^2 = \frac{1}{L}
  \sum_{l=1}^{L} \|\mathbf{R}\bar{\mathbf{p}}_l - \bar{\mathbf{q}}_l\|^2
\label{eq:rmsd}
\end{equation}
is given by
\begin{equation}
\mathbf{R}^* = \mathbf{V}
  \begin{pmatrix} 1 & & \\ & 1 & \\ & & \det(\mathbf{V}\mathbf{W}^T) \end{pmatrix}
  \mathbf{W}^T,
\label{eq:kabsch-rotation}
\end{equation}
where $\mathbf{H} = \bar{\mathbf{P}}^T \bar{\mathbf{Q}} \in \mathbb{R}^{3 \times 3}$
is the cross-covariance matrix and $\mathbf{H} = \mathbf{W}\boldsymbol{\Sigma}\mathbf{V}^T$
is its SVD.
\end{theorem}

\begin{proof}
Standard; see Kabsch (1976).
The key observation is that $\mathbf{H}$ is always $3 \times 3$
regardless of $L$, so its SVD costs $O(1)$.
The total cost is dominated by forming $\mathbf{H}$, which
requires $O(L)$ vector operations.
\end{proof}

\begin{corollary}[Stage 2 Complexity]
\label{cor:stage2}
Given the atom map $\mathcal{M}$ from Stage~1, the 3D
alignment and RMSD computation cost $O(L) = O(n)$.
\end{corollary}

\begin{definition}[Hybrid Kernel]
\label{def:hybrid}
The two-stage similarity measure combines 2D topological
and 3D geometric information:
\begin{equation}
k_{\text{hybrid}}(G, H)
  = (1 - \beta)\, k_{\text{2D}}(G, H)
  + \beta\, k_{\text{3D}}(G, H),
\label{eq:hybrid-kernel}
\end{equation}
where
$k_{\text{2D}} = \exp(-\gamma_{\text{2D}} F^*/L)$ is the
topological kernel from~\eqref{eq:kernel},
$k_{\text{3D}} = \exp(-\gamma_{\text{3D}} \cdot \mathrm{RMSD})$
is the geometric kernel,
and $\beta \in [0,1]$ is cross-validated per target.
\end{definition}

\begin{remark}[Overall Pipeline Complexity]
The total per-pair cost of the hierarchical pipeline is
\begin{equation}
\underbrace{O(T_\zeta T_{\mathrm{FW}} n^3)}_{\text{Stage 1: 2D kernel}}
+ \underbrace{O(n)}_{\text{Stage 2: 3D alignment}}
= O(T_\zeta T_{\mathrm{FW}} n^3),
\label{eq:pipeline-complexity}
\end{equation}
completely dominated by Stage~1.
The 2D stage serves as a search-space reduction gate:
only candidates with $k_{\text{2D}} > \tau$
(a tunable threshold) proceed to Stage~2.
For a library of $P$ molecules with filtration rate~$\phi$,
the total screening cost is
\begin{equation}
\underbrace{O\!\left(\binom{P}{2} T_\zeta T_{\mathrm{FW}} n^3\right)}_{\text{2D pairwise kernel matrix}}
+ \underbrace{O(\phi P^2 n)}_{\text{3D alignment of filtered pairs}}.
\label{eq:screening-cost}
\end{equation}
With aggressive filtering ($\phi \!\ll\! 1$), the 3D
cost becomes negligible.
\end{remark}


% ===================================================================
\section{Relabeling Invariance}
\label{sec:invariance}
% ===================================================================

\begin{proposition}[Relabeling Invariance]
\label{prop:invariance}
Let $\mathbf{P} \!\in\! \{0,1\}^{M \times M}$ and
$\mathbf{Q} \!\in\! \{0,1\}^{N \times N}$ be permutation
matrices, and let
$\mathbf{A}'_G = \mathbf{P}\mathbf{A}_G\mathbf{P}^T$,
$\mathbf{A}'_H = \mathbf{Q}\mathbf{A}_H\mathbf{Q}^T$.
Then for any $\mathbf{X} \!\in\! \mathcal{D}_L$,
\begin{equation}
F\!\bigl(\mathbf{P}\mathbf{X}\mathbf{Q}^T;\,
  \mathbf{A}'_G, \mathbf{A}'_H, \mathbf{C}'\bigr)
= F\!\bigl(\mathbf{X};\,
  \mathbf{A}_G, \mathbf{A}_H, \mathbf{C}\bigr),
\label{eq:invariance}
\end{equation}
where $\mathbf{C}' = \mathbf{P}\mathbf{C}\mathbf{Q}^T$.
\end{proposition}

\begin{proof}
Let $\mathbf{X}' = \mathbf{P}\mathbf{X}\mathbf{Q}^T$.
Then $\mathbf{X}'\mathbf{A}'_H(\mathbf{X}')^T
= \mathbf{P}\mathbf{X}\mathbf{A}_H\mathbf{X}^T\mathbf{P}^T$
and $\mathbf{U}' \circ \mathbf{A}'_G
= \mathbf{P}(\mathbf{U}\!\circ\!\mathbf{A}_G)\mathbf{P}^T$.
The Frobenius norm is unitarily invariant:
$\|\mathbf{P}\mathbf{M}\mathbf{P}^T\|_F = \|\mathbf{M}\|_F$.
The trace term satisfies
$\operatorname{tr}((\mathbf{C}')^T\mathbf{X}')
= \operatorname{tr}(\mathbf{C}^T\mathbf{X})$
by the cyclic property of trace.
\end{proof}


% ===================================================================
\section{Full Algorithm}
\label{sec:algorithm}
% ===================================================================

{\small
\begin{algorithm}[H]
\caption{WCS-GNCCP with Frank-Wolfe}
\label{alg:gnccp-fw}
\begin{algorithmic}[1]
\REQUIRE $\mathbf{A}_G,\mathbf{A}_H$; cost $\mathbf{C}$;
  size $L$; weights $\alpha,\gamma$
\ENSURE Similarity score $k \!\in\! (0, 1]$
  and atom map $\mathcal{M}$
\STATE $\hat{\mathbf{A}}_G \leftarrow
  \mathbf{A}_G / \|\mathbf{A}_G\|_F$;\;
  $\hat{\mathbf{A}}_H \leftarrow
  \mathbf{A}_H / \|\mathbf{A}_H\|_F$
  \hfill\COMMENT{Prop.~\ref{prop:normalization}}
\STATE $\hat{\mathbf{C}} \leftarrow
  (\mathbf{C}\!-\!C_{\min})/(C_{\max}\!-\!C_{\min})$
\STATE $\mathbf{X} \leftarrow
  \frac{L}{MN}\mathbf{1}_{M\times N}$
  \hfill\COMMENT{Prop.~\ref{prop:endpoints}(1)}
\FOR{$\zeta = 1, 1\!-\!\Delta\zeta, \ldots, -1$}
  \FOR{$t = 1, \ldots, T_{\text{FW}}$}
    \STATE $\nabla \leftarrow
      \nabla_{\!\mathbf{X}} J_\zeta$
      \hfill\COMMENT{Prop.~\ref{prop:grad-Jzeta}}
    \STATE $\mathbf{Y} \leftarrow
      \text{Hungarian}(\nabla)$
      \hfill\COMMENT{Def.~\ref{def:lmo}}
    \STATE $g \leftarrow
      \langle \nabla,\,\mathbf{X}\!-\!\mathbf{Y}\rangle$
      \hfill\COMMENT{Thm.~\ref{thm:fw-gap}}
    \IF{$g < \varepsilon$}
      \STATE \textbf{break}
    \ENDIF
    \STATE $\lambda \leftarrow
      \text{ArmijoLS}(\mathbf{X},\mathbf{Y},J_\zeta)$
    \STATE $\mathbf{X} \leftarrow
      \text{Sinkhorn}\!\bigl(
        \mathbf{X}\!+\!\lambda(\mathbf{Y}\!-\!\mathbf{X}),\,L
      \bigr)$
      \hfill\COMMENT{Prop.~\ref{prop:sinkhorn}}
  \ENDFOR
\ENDFOR
\STATE $\mathbf{X}^{\!*} \leftarrow
  \text{Hungarian}(-\mathbf{X})$
  \hfill\COMMENT{discretize}
\STATE $\mathcal{M} \leftarrow
  \{(i,j) : X^*_{ij} = 1\}$
  \hfill\COMMENT{atom map}
\STATE $F^{\!*} \leftarrow
  F(\mathbf{X}^{\!*})$
  \hfill\COMMENT{Def.~\ref{def:wcs-objective}}
\RETURN $k = \exp(-\gamma F^{\!*}/L)$,\;$\mathcal{M}$
  \hfill\COMMENT{Prop.~\ref{prop:psd}}
\end{algorithmic}
\end{algorithm}
}


\begin{thebibliography}{9}
\bibitem{schoenberg1938metric}
I.~J.~Schoenberg,
``Metric spaces and positive definite functions,''
\textit{Trans.\ Amer.\ Math.\ Soc.},
vol.~44, no.~3, pp.~522--536, 1938.

\bibitem{kabsch1976solution}
W.~Kabsch,
``A solution for the best rotation to relate two sets of vectors,''
\textit{Acta Crystallographica Section A},
vol.~32, no.~5, pp.~922--923, 1976.

\bibitem{frank1956algorithm}
M.~Frank and P.~Wolfe,
``An algorithm for quadratic programming,''
\textit{Naval Research Logistics Quarterly},
vol.~3, no.~1--2, pp.~95--110, 1956.

\bibitem{liu2015graduated}
Z.-Y.~Liu and H.~Qiao,
``GNCCP --- graduated non-convexity and concavity procedure,''
\textit{IEEE Trans.\ Pattern Anal.\ Mach.\ Intell.},
vol.~36, no.~6, pp.~1258--1267, 2014.

\bibitem{kuhn1955hungarian}
H.~W.~Kuhn,
``The Hungarian method for the assignment problem,''
\textit{Naval Research Logistics Quarterly},
vol.~2, no.~1--2, pp.~83--97, 1955.
\end{thebibliography}

\end{document}
